\documentclass[11pt,a4paper]{report}

\usepackage[utf8]{inputenc}
\usepackage[T1]{fontenc}
\usepackage[french]{babel}

\usepackage{newtxtext,newtxmath} 

\usepackage{geometry}
\geometry{
    paper=a4paper,
    top=2.5cm,
    bottom=2.5cm,
    left=2.5cm,
    right=2.5cm,
    headheight=15pt, 
    footskip=30pt
}

\usepackage{setspace}
\onehalfspacing 

% --- Gestion des tailles de titres ---
\usepackage{titlesec}
\titleformat{\chapter}[block]
  {\normalfont\fontsize{14}{16}\bfseries}{\thechapter.}{0.5em}{}
\titlespacing*{\chapter}{0pt}{-20pt}{20pt}

\titleformat{\section}
  {\normalfont\fontsize{12}{14}\bfseries}{\thesection}{1em}{}

\usepackage{fancyhdr}
\pagestyle{fancy}
\fancyhf{} 

\fancyhead[L]{\fontsize{10}{12}\selectfont Licence 3 Informatique – Université de Lille}
\fancyfoot[R]{\thepage}
\fancyfoot[L]{\fontsize{10}{12}\selectfont Stella-Rose MILLE} 

\fancypagestyle{plain}{
    \fancyhf{}
    \fancyhead[L]{\fontsize{10}{12}\selectfont Licence 3 Informatique – Université de Lille}
    \fancyfoot[R]{\thepage}
    \fancyfoot[L]{\fontsize{10}{12}\selectfont Stella-Rose MILLE}
    \renewcommand{\headrulewidth}{0.4pt}
}
\renewcommand{\headrulewidth}{0.4pt}

% --- Figures et Légendes ---
\usepackage{graphicx}
\usepackage{caption}
\captionsetup{font=small, labelfont=bf}

\usepackage[colorlinks=true, linkcolor=black, citecolor=blue, urlcolor=blue]{hyperref}


\begin{document}

% ======================================================================
% PAGE DE GARDE
% ======================================================================
\begin{titlepage}
    \centering
    
    \begin{minipage}{0.45\textwidth}
        \begin{flushleft}
            \vspace{0.2cm}
            \fontsize{9}{11}\selectfont
            \textbf{Université de Lille}\\
            Faculté des Sciences et Technologies\\
            Cité Scientifique, Avenue Henri Poincaré\\
            59655 Villeneuve-d'Ascq, France
        \end{flushleft}
    \end{minipage}
    \hfill
    \begin{minipage}{0.45\textwidth}
        \begin{flushright}
            \vspace{0.2cm}
            \fontsize{9}{11}\selectfont
            \textbf{INRIA Lille -- Nord Europe}\\
            Parc Scientifique de la Haute Borne\\
            Équipe EVREF\\
        \end{flushright}
    \end{minipage}
    
    \vfill
    \vspace{0.5cm}
    
    % --- Intitulé Précis de la Formation ---
    {\fontsize{14}{16}\selectfont\bfseries Licence 3 Informatique}\\[0.4cm]
    {\fontsize{11}{13}\selectfont Année Universitaire 2025 -- 2026}\\[1.5cm]
    
    % --- Ligne de séparation ---
    \rule{\textwidth}{1pt}\\[0.4cm]
    % --- Titre du Rapport ---
    {\fontsize{13}{15}\selectfont\scshape RAPPORT DE STAGE INTERMÉDIAIRE}\\[0.5cm]
    {\fontsize{18}{22}\selectfont\bfseries Évolution et optimisation d'un framework de Property-Based Testing pour Pharo}\\[0.6cm]
    \rule{\textwidth}{1pt}\\[0.8cm]
    
    % --- Période et Dates du Stage ---
    {\fontsize{12}{14}\selectfont\bfseries Période du stage :}\\
    {\fontsize{11}{13}\selectfont du 13 avril au 28 août 2026}\\[2cm]
    
    % --- Bloc Acteurs (Auteur et Tuteurs) ---
    \begin{minipage}{0.48\textwidth}
        \begin{flushleft} \large
            \emph{Rédigé et présenté par :}\\
            \textbf{Stella-Rose MILLE}\\[0.2cm]
        \end{flushleft}
    \end{minipage}
    \hfill
    \begin{minipage}{0.48\textwidth}
        \begin{flushright} \large
            \emph{Tuteur en Entreprise :}\\
            \textbf{Federico LOCHBAUM}\\[0.6cm]
            
            \emph{Tuteur Universitaire :}\\
            \textbf{Romain ROUVOY}
        \end{flushright}
    \end{minipage}
    
    \vfill
\end{titlepage}

% ======================================================================
% TABLE DES MATIERES
% ======================================================================
\tableofcontents

% ======================================================================
% CONTEXTE DU STAGE 
% ======================================================================
\chapter{Contexte du Stage}
\label{ch:contexte}

\section{Présentation du Laboratoire et de l'Équipe}
Ce stage se déroule au sein d'INRIA (Institut National de Recherche en Sciences et Technologies du Numérique), un organisme public de recherche dédié aux sciences du numérique. 
Plus précisément, j'ai intégré l'équipe de recherche Evref. 

Leur principal objectif est l'étude de l'évolution des grands systèmes 
informatiques. Leurs travaux s'articulent autour de trois axes :
L'analyse et la migration, le développement d'outils, 
l'infrastructure.

\section{Contexte de l'Application : Le projet Ume}

Mon travail s'inscrit dans le cadre d'un projet de thèse de l'équipe, 
centré sur le développement d'Ume. Ume est un framework de Property-Based Testing
 (PBT) conçu spécifiquement pour l'écosystème Pharo 
 (un langage purement orienté objet dérivé de Smalltalk). 
 Contrairement aux tests unitaires classiques où le développeur 
 écrit manuellement les entrées, le PBT génère automatiquement des 
 centaines d'entrées aléatoires pour valider des propriétés globales du logiciel. 
 Ume cherche à pousser ce concept plus loin en l'appliquant à la détection 
 de régressions et d'anomalies de performance.

% ======================================================================
% DESCRIPTION D'UNE REALISATIONS
% ======================================================================
\chapter{Description d'une Réalisation Technique Majeure}
\label{ch:realisation}
Au cours de ces sept premières semaines de stage, j'ai réalisé plusieurs 
améliorations centrées sur l'API d'Ume afin de simplifier son adoption 
par des utilisateurs extérieurs. J'ai notamment fluidifié la définition 
d'une campagne de test et implémenté une barre de progression 
multithreadée pour offrir un retour visuel en temps réel lors de l'exécution. 
Cependant, ma réalisation majeure réside dans l'intégration du mécanisme de 
shrinking.

\section{Contexte et Objectif du Shrinking}
Le shrinking intervient lorsqu'un test échoue sur une entrée complexe 
générée aléatoirement. L'objectif est de réduire cet input au cas minimal 
reproduisant l'erreur. Bien que l'application du shrinking aux problématiques 
pures de performance reste un sujet de recherche ouvert, cette 
fonctionnalité est indispensable pour permettre à l'utilisateur de 
localiser l'origine précise d'un bug fonctionnel sans être noyé sous des données 
superflues.

\section{Choix Techniques et Justifications}
L'implémentation du shrinking dans Ume s'inspire des travaux de référence du
Fuzzing Book. 
Mon premier choix s'est porté sur la création d'un algorithme de 
Delta Debugging. Basé sur le principe du « diviser pour régner », 
cet algorithme segmente la chaîne de caractères d'entrée et évalue chaque morceau 
pour isoler la faille. 

Cependant, le Delta Debugging classique montre ses limites lorsqu'il doit 
manipuler des entrées structurées. Si la réduction casse la syntaxe attendue 
par le programme, le test échoue pour une mauvaise raison (erreur de parsing). 
Pour résoudre cette problématique, nous avons choisi de coupler cette approche 
à des \textbf{réducteurs basés sur la grammaire}, garantissant que chaque entrée 
réduite reste syntaxiquement valide.

\section{Développements et Travaux Réalisés}
Pour mener à bien cette réalisation, j'ai adopté une méthodologie rigoureuse 
de développement piloté par les tests (TDD). Cela m'a permis de 
définir précisément le comportement attendu de mes réducteurs avant d'écrire 
le code métier. J'ai mis en place le design pattern Strategy pour 
l'architecture du shrinker. Ce choix permet d'isoler les différents 
algorithmes de réduction (Delta Debugging, réduction de grammaire) au sein 
de classes distinctes, rendant le framework facilement extensible pour de 
futurs algorithmes.

\section{Difficultés rencontrées et Solutions apportées}
\begin{itemize}
    \item \textbf{Manque de documentation et de détails théoriques :} Les 
    articles académiques détaillent rarement l'implémentation bas niveau des 
    algorithmes. J'ai dû analyser des codebases open-source existants pour en 
    comprendre la logique fine.
    \item \textbf{Adaptation au paradigme Smalltalk/Pharo :} Les exemples de 
    référence (notamment du \textit{Fuzzing Book}) sont écrits en Python. 
    Transposer une logique impérative ou procédurale vers le modèle purement 
    orienté objet de Pharo a demandé une réingénierie complète des structures 
    de données.
    \item \textbf{Solidité de l'exploration :} Garantir que le shrinker explore 
    efficacement les branches de réduction sans boucler indéfiniment a nécessité 
    la mise en place d'une suite de tests aux limites très stricte.
\end{itemize}

% ======================================================================
% CONCLUSION ET BILAN 
% ======================================================================
\chapter{Conclusion et Bilan du Stage}
\label{ch:conclusion}

\section{Adéquation de la Réalisation aux Objectifs et Perspectives}
Les objectifs de cette première moitié de stage ont été atteints : l'API 
d'Ume a été simplifiée pour l'utilisateur final et une première brique 
fonctionnelle de shrinking respectant la grammaire a été intégrée. Concernant 
la question de l'évaluation, la mise en place de \textbf{benchmarks} précis 
constitue l'étape logique suivante. Ils permettront de mesurer l'impact du 
shrinking sur le temps de calcul et le nombre de tests générés.

\section{Présentation du Travail à Réaliser pour la Suite}
La suite de mes missions s'articulera autour de deux axes :
\begin{itemize}
    \item L'implémentation complète et l'optimisation des algorithmes de 
    réduction de grammaire existants.
    \item La réalisation de campagnes de Benchmarks pour quantifier les 
    performances du shrinker.
    \item L'objectif à terme est de concevoir un système capable d'aiguiller 
    intelligemment la campagne de shrinking vers l'algorithme le plus adapté 
    au cas par cas, afin de limiter le coût computationnel actuel de cette 
    procédure.
\end{itemize}

\section{Enseignements et Apports du Stage}
Ce stage m'a permis de développer une forte autonomie. Dès les premiers jours, 
mon tuteur a mis l'accent sur l'auto-organisation. J'ai appris à chercher 
l'information par moi-même en lisant des articles de recherche. 
Nous avons structuré notre flux de travail avec GitHub pour la gestion de 
versions (apprentissage des Pull Requests professionnelles et bien documentées) 
et \textbf{Obsidian} pour la documentation de mes tâches quotidiennes. 

Sur le plan technique, l'utilisation rigoureuse du TDD m'a permis de canaliser 
mes développements lorsque je perdais de vue l'objectif final, tandis que 
l'application concrète des Design Patterns a renforcé mes compétences en 
conception logicielle.

\section{Remerciements}
Je souhaite remercier mon tuteur d'entreprise, Federico Lochbaum, pour ses 
précieux conseils, sa patience et sa confiance. Mon tuteur universitaire, 
Romain Rouvoy, pour son suivi rigoureux. Enfin, j'adresse mes remerciements 
à Imen Sayar pour son soutien continu et pour avoir appuyé ma candidature 
pour ce stage.

\end{document}